%% Results in Engineering (Elsevier) -- elsarticle template
\documentclass[preprint,12pt,authoryear]{elsarticle}

\usepackage[utf8]{inputenc}
\usepackage[T1]{fontenc}
\usepackage{amsmath,amssymb,amsfonts,bm}
\usepackage{graphicx}
\usepackage{booktabs}
\usepackage{longtable}
\usepackage{multirow}
\usepackage{xcolor}
\usepackage{tikz}
\usetikzlibrary{positioning,arrows.meta,fit,backgrounds,calc,shapes.geometric}
\usepackage[hidelinks]{hyperref}
\usepackage[margin=2.3cm]{geometry}

\graphicspath{{figures/}}
\newcommand{\inputtable}[1]{\IfFileExists{tables/#1.tex}{\input{tables/#1.tex}}{\begin{table}[h]\centering\fbox{\parbox{0.8\linewidth}{\centering Table \texttt{\detokenize{#1}} is generated by \texttt{scripts/05\_report.py}.}}\end{table}}}
\newcommand{\fig}[2][\linewidth]{\IfFileExists{figures/#2.pdf}{\includegraphics[width=#1]{#2}}{\fbox{\parbox{0.8\linewidth}{\centering Figure \texttt{\detokenize{#2}} is generated by \texttt{scripts/05\_report.py}.}}}}
\newcommand{\ind}{\mathbf{1}}
\newcommand{\dA}{\delta_A}
\newcommand{\dB}{\delta_B}

\setlength{\emergencystretch}{2em}
\journal{Results in Engineering}

\begin{document}

\begin{frontmatter}

\title{Where are the structures and how do they evolve? Pre-processing, segmentation architectures and post-processing for the Gray--Scott reaction--diffusion system of \emph{The Well}}

\author[ufma]{Wesley dos Santos Gatinho}
\affiliation[ufma]{organization={Universidade Federal do Maranh\~ao (UFMA)},
            addressline={Av. dos Portugueses, 1966, Bacanga},
            city={S\~ao Lu\'is},
            postcode={65080-805},
            state={MA},
            country={Brazil}}

\begin{abstract}
Reaction--diffusion systems underlie models of chemical reactors, of the spreading of reacting substances and of biological pattern formation, and quantifying \emph{where} their coherent structures are and \emph{how} they evolve requires segmenting the concentration fields over time. Following a protocol shared by a class, we evaluate how pre-processing, segmentation architecture and post-processing affect the segmentation of the active-reaction regions of the two-dimensional Gray--Scott system of \emph{The Well}. A physically defined ground truth thresholds the autocatalytic species $B$ with Otsu's method at a characteristic mid-simulation frame, and the networks recover it from the substrate $A$ at its native $128\times128$ resolution. Ten pre-processing methods, three architectures (U-Net, SegFormer and Mask2Former, 20 training epochs each) and five post-processing methods were applied to all 1,200 trajectories, giving the 150 combinations of the protocol matrix plus baseline rows. On 120 test trajectories, U-Net reached a Dice of 0.953 without pre- or post-processing, against 0.892 for SegFormer, 0.768 for Mask2Former and 0.888 for Otsu thresholding of $A$. Pre-processing changed Dice by at most 1.37 percentage points, temporal filtering lowered it by 8.6 to 11.3 points, and the largest gains came from the random initial snapshot: excluding it, threshold optimization lowered the Dice of U-Net and SegFormer. Selection of configurations by mean Dice was therefore sensitive to a single frame without physical information.
\end{abstract}

\begin{keyword}
Gray--Scott model \sep reaction--diffusion \sep Turing patterns \sep semantic segmentation \sep U-Net \sep SegFormer \sep Mask2Former \sep The Well
\end{keyword}

\end{frontmatter}

%% ===================================================================
\section{Introduction}
\label{sec:intro}

Since Turing's seminal work on morphogenesis \citep{turing1952}, reaction--diffusion (RD) equations have been the canonical framework to explain how spatial order emerges from an almost homogeneous medium. Diffusion, which alone smooths every perturbation, becomes destabilising when it is coupled to a nonlinear local reaction and the species diffuse at different rates. The result is a variety of stationary and dynamic patterns---spots, stripes, labyrinths, self-replicating spots, travelling ``gliders'' and spiral waves---used to explain biological patterns \citep{kondo2010,murray2003} and observed in chemical reactors \citep{lee1993}. The Gray--Scott system \citep{gray1983,gray1984} is the reference minimal model of this phenomenon: two species, one cubic autocatalytic reaction and two control parameters generate the patterns catalogued by \citet{pearson1993}. Analysing such a pattern requires answering where the structures of interest are and how they evolve over time, i.e.\ turning the continuous field into regions by segmentation.

\emph{The Well} \citep{ohana2024well} provides Gray--Scott simulations for machine learning, where they are used to benchmark neural surrogates such as the Fourier Neural Operator \citep{li2021fno}, i.e.\ to \emph{predict} future fields \citep{brunton2020}. Prediction does not locate the structures nor measure their formation, migration and extinction. Deep segmentation networks, established in natural and medical images \citep{long2015fcn,ronneberger2015unet,xie2021segformer,cheng2022mask2former}, provide that representation, but their behaviour on simulated physical fields---smooth, textureless, periodic and with structures that appear and vanish---is less documented. Moreover, classical image-processing steps that precede or follow the network are often applied without measuring whether and why they change the segmentation.

This work evaluates the effect of pre-processing, segmentation architecture and post-processing on the temporal segmentation of the active-reaction regions of the Gray--Scott system of The Well. It follows a protocol shared by a Digital Image Processing class, in which every participant applies the same experiment to a different 2D system: ten pre-processing methods, three architectures and five post-processing methods; an input of at most 128 pixels along the largest axis, without upscaling; a subset of at most 25~GB per dataset; and 20 training epochs. The experiment has four conditions: the architecture alone, pre-processing + architecture, architecture + post-processing, and pre-processing + architecture + post-processing. Four research questions guide the study: (RQ1) how U-Net, SegFormer and Mask2Former compare with each other and with a learning-free reference; (RQ2) to what extent pre-processing changes the segmentation; (RQ3) to what extent post-processing changes it; and (RQ4) whether the masks follow the temporal evolution of the structures, including their formation and extinction.

The contribution is empirical and descriptive: a physically defined ground truth, the evaluation of the 150 combinations of the protocol and of their baseline rows with per-class overlap metrics, temporal metrics adapted to periodic domains and computational cost, paired analyses that isolate the effect of each technique, and a sensitivity analysis to the initial snapshot of the simulation. No new architecture is proposed. The results are delimited by one random seed per training, by a threshold-based rather than expert-annotated ground truth, and by the adaptation of the Mask2Former training to the available budget (Section~\ref{sec:methods}).

%% ===================================================================
\section{Where the Gray--Scott system applies}
\label{sec:applications}

\paragraph{Chemical reaction engineering} The model was derived for an isothermal continuously stirred tank reactor with cubic autocatalysis, $A+2B\rightarrow 3B$, $B\rightarrow P$ \citep{gray1983,gray1984}, to explain multistability, ignition, extinction and oscillations. Its spatially extended form describes gel reactors, where the self-replicating spots predicted by the model \citep{pearson1993} were observed in the ferrocyanide--iodate--sulfite reaction \citep{lee1993}. The segmented region $\{B>\tau\}$ is the region where the reaction is ignited; its area over time and its possible collapse relate to reactor yield and operation.

\paragraph{Industrial safety} After an accidental release, the concentration of a toxic gas spreads by diffusion while it reacts with the air, following equations with the same diffusion + reaction structure (plus advection under wind). Emergency planning simulates such scenarios---for example an ammonia leak in a plant---to define evacuation zones and place sensors, which requires delimiting, at every instant, the region where a simulated concentration exceeds a threshold and tracking how it grows and fades.

\paragraph{Environmental disasters and biology} Harmful algal blooms grow by consuming dissolved nutrients and spread in the water, a consumer--resource coupling analogous to the pair $A$ and $B$. Turing-type models explain pigmentation patterns and the spacing of follicles \citep{kondo2010,murray2003}; segmentation turns such patterns into measurable quantities.

\paragraph{Scientific machine learning} The Well uses this system to benchmark surrogates \citep{ohana2024well}; a reliable segmentation provides a structure-level diagnostic of whether patterns are preserved in long rollouts.

%% ===================================================================
\section{Theoretical background}
\label{sec:theory}

\subsection{Reaction--diffusion and the Turing mechanism}

A two-species RD system reads
\begin{equation}
\frac{\partial u}{\partial t} = D_u \nabla^2 u + f(u,v), \qquad
\frac{\partial v}{\partial t} = D_v \nabla^2 v + g(u,v),
\end{equation}
where $\nabla^2$ is the Laplacian (Fickian diffusion) and $f,g$ are the local kinetics. Let $\mathbf{J}$ be the Jacobian of the kinetics at a homogeneous steady state. A perturbation $\propto e^{\lambda t + i\mathbf{q}\cdot\mathbf{x}}$ grows with the eigenvalues of $\mathbf{J}-q^2\mathbf{D}$, $\mathbf{D}=\mathrm{diag}(D_u,D_v)$. A state stable without diffusion ($\operatorname{tr}\mathbf{J}<0$, $\det\mathbf{J}>0$) is destabilised by diffusion if
\begin{equation}
h(q^2)=D_uD_v q^4-(D_u J_{22} + D_v J_{11})\,q^2+\det\mathbf{J} < 0 \quad\text{for some } q,
\end{equation}
which requires sufficiently different diffusivities \citep{turing1952}. The band of unstable wavenumbers sets the characteristic length of the pattern \citep{murray2003}.

\subsection{The Gray--Scott model}

A substrate $A$ is converted by the autocatalytic step $A+2B\rightarrow3B$ and the autocatalyst decays as $B\rightarrow P$ \citep{gray1983,gray1984}:
\begin{align}
\frac{\partial A}{\partial t} &= \dA\nabla^2 A - A B^2 + F\,(1-A), \label{eq:gsA}\\
\frac{\partial B}{\partial t} &= \dB\nabla^2 B + A B^2 - (F+k)\,B, \label{eq:gsB}
\end{align}
where $AB^2$ is the reaction rate, $F$ the feed rate of $A$ and $k$ the kill rate of $B$. The trivial state $(1,0)$ always exists and is stable (eigenvalues $-F$ and $-(F+k)$). Non-trivial states satisfy $AB=F+k$ and $F(1-A)=(F+k)B$,
\begin{equation}
B_\pm=\frac{F\pm\sqrt{F^2-4F(F+k)^2}}{2(F+k)},\qquad A_\pm=\frac{F+k}{B_\pm},
\end{equation}
and exist below the saddle--node curve $k=\tfrac12\sqrt{F}-F$, with Jacobian
\begin{equation}
\mathbf{J}=\begin{pmatrix} -B^{*2}-F & -2(F+k)\\ B^{*2} & F+k\end{pmatrix}.
\end{equation}
Close to these boundaries, localised perturbations evolve into self-replicating spots, travelling pulses and labyrinths depending on $(F,k)$ \citep{pearson1993}; for some parameters the pattern collapses to the trivial state and all structures disappear. Because the reaction consumes $A$ where $B$ is high, the fields are anti-correlated; because $A$ diffuses twice as fast as $B$, the depletion halo of $A$ is wider and smoother than the structure, so the boundary $\{B=\tau\}$ is not an iso-line of $A$ and recovering it from $A$ is a non-linear, non-local inverse problem.

\subsection{The \texttt{gray\_scott\_reaction\_diffusion} dataset}

The dataset has 1,200 trajectories: six parameter sets (regimes), each with 200 random initial conditions, split by trajectory into 160/20/20 \citep{ohana2024well}: \emph{gliders} ($F=0.014,k=0.054$), \emph{bubbles} ($0.098,0.057$), \emph{maze} ($0.029,0.057$), \emph{worms} ($0.058,0.065$), \emph{spirals} ($0.018,0.051$) and \emph{spots} ($0.030,0.062$). All use $\dA=2\times10^{-5}$, $\dB=10^{-5}$ on the periodic square $[-1,1]^2$ with $128\times128$ points, a Fourier pseudo-spectral discretisation and ETDRK4 time stepping \citep{cox2002,kassam2005}. Snapshots are stored every 10 time units up to $t=10{,}000$ (1,001 snapshots per trajectory, 154~GB); we index time by the snapshot number $n$.

\subsection{Segmentation architectures}

\paragraph{U-Net} A symmetric encoder--decoder whose skip connections re-inject full-resolution detail \citep{ronneberger2015unet}; here with a ResNet-34 encoder pre-trained on ImageNet \citep{he2016resnet,deng2009imagenet} (residual blocks $\mathbf{y}=\mathcal{F}(\mathbf{x})+\mathbf{x}$; 24.4~M parameters) and a full-resolution logit map, $p=\sigma(z)$.

\paragraph{SegFormer} Self-attention, $\mathrm{softmax}(QK^\top/\sqrt{d})V$ \citep{vaswani2017,dosovitskiy2021vit}, in a hierarchical Mix Transformer encoder \citep{wang2021pvt,xie2021segformer} with overlapping patch merging, efficient self-attention with spatial reduction $R$ (cost $O(N^2/R)$), Mix-FFN, $\mathbf{x}_{out}=\mathrm{MLP}(\mathrm{GELU}(\mathrm{Conv}_{3\times3}(\mathrm{MLP}(\mathbf{x}_{in}))))+\mathbf{x}_{in}$, and an all-MLP decoder that predicts at $1/4$ resolution. MiT-B2, ImageNet-1k, 27.3~M parameters.

\paragraph{Mask2Former} Mask classification: $N$ pairs $\{(p_i,m_i)\}$ of class distributions (with a ``no object'' class) and masks \citep{cheng2021maskformer,carion2020detr}, with a Swin-T backbone \citep{liu2021swin}, a multi-scale deformable-attention pixel decoder \citep{zhu2021deformable} and a Transformer decoder with masked attention \citep{cheng2022mask2former},
\begin{equation}
\mathbf{X}_l=\mathrm{softmax}(\mathcal{M}_{l-1}+\mathbf{Q}_l\mathbf{K}_l^\top)\mathbf{V}_l+\mathbf{X}_{l-1},\quad
\mathcal{M}_{l-1}(x,y)=\begin{cases}0 & m_{l-1}(x,y)\geq 0.5\\ -\infty & \text{otherwise.}\end{cases}
\end{equation}
Predictions are matched to ground-truth segments by the Hungarian algorithm \citep{kuhn1955}; the loss is $\lambda_{ce}\mathcal{L}_{ce}+\lambda_{bce}\mathcal{L}_{bce}+\lambda_{dice}\mathcal{L}_{dice}$ ($2,5,5$), applied in the original recipe also to every intermediate decoder layer (deep supervision). The semantic score of class $c$ is $\sum_i p_i(c)m_i$. Swin-T from ADE20K \citep{zhou2017ade}, 47.4~M parameters, $N=100$, two classes.

\paragraph{Output resolution} U-Net predicts at the input resolution, whereas SegFormer and Mask2Former predict at $1/4$ of it ($32\times32$ for a $128\times128$ input), bilinearly resized to $128\times128$.

\subsection{Pre-processing and post-processing}

Value normalisations use training statistics: z-score $x'=(x-\mu)/\sigma$; RMS normalisation $x'=x/\sqrt{\overline{x^2}}$, which preserves zero and the shape of the distribution; robust scaling $x'=(x-\mathrm{med})/\mathrm{IQR}$; min--max scaling; and percentile clipping to $[p_1,p_{99}]$. Spatial filters \citep{gonzalez2018}: Gaussian smoothing (spectrum multiplied by $e^{-\sigma^2q^2/2}$); median filtering; bilateral filtering, which averages pixels that are both close and similar \citep{tomasi1998}; Perona--Malik anisotropic diffusion, $\partial_\tau I=\nabla\cdot(c(|\nabla I|)\nabla I)$, $c(s)=e^{-(s/\kappa)^2}$ \citep{perona1990}; and wavelet denoising with soft thresholding $\eta_\lambda(w)=\mathrm{sign}(w)\max(|w|-\lambda,0)$ and the BayesShrink threshold $\lambda=\hat\sigma_n^2/\hat\sigma_x$ \citep{donoho1994,chang2000,daubechies1992}.

Post-processing turns the probability map $P$ into the final mask: threshold optimisation chooses $\tau$ in $\ind[P>\tau]$; morphological opening $\gamma_B(X)=(X\ominus B)\oplus B$ removes objects smaller than $B$, and closing $\varphi_B(X)=(X\oplus B)\ominus B$ fills narrow gaps \citep{serra1982,soille2003}; connected-component filtering removes regions below a minimum area; and temporal consistency filtering takes the median of the probabilities over $w$ consecutive snapshots of a trajectory before thresholding.

\subsection{Evaluation metrics}
\label{sec:metrics}

With $TP,FP,FN,TN$ the pixel counts of a frame, the main metrics for the structure class are $\mathrm{IoU}=TP/(TP+FP+FN)$ \citep{jaccard1912} and $\mathrm{Dice}=2TP/(2TP+FP+FN)$ \citep{dice1945}, related by $\mathrm{Dice}=2\,\mathrm{IoU}/(1+\mathrm{IoU})$ \citep{taha2015,maierhein2024}. We also report precision, recall, the background-class metrics $\mathrm{IoU}_{bg}=TN/(TN+FP+FN)$ and $\mathrm{Dice}_{bg}$, and mIoU, the mean IoU of the two classes. Frames where prediction and ground truth are both empty count as perfect. Scores are averaged over frames of a trajectory and then over trajectories.

Temporal metrics compare the dynamics of predicted and true masks: the area-variation error $\mathrm{AVE}=\frac1T\sum_j|\Delta a_{pred}(t_j)-\Delta a_{gt}(t_j)|$ ($a$ in \% of the domain); the structure-centroid displacement error (CDE), which merges components across periodic borders, computes circular-mean centroids $\bar y=\frac{H}{2\pi}\operatorname{atan2}(\sum\sin\theta_y,\sum\cos\theta_y)$ and pairs each true and predicted centroid with its nearest counterpart using the periodic distance $d_{per}$,
\begin{equation}
\mathrm{CDE}(t)=\tfrac12\Big[\tfrac{1}{|\mathcal{G}|}\textstyle\sum_{g}\min_{p}d_{per}(g,p)+\tfrac{1}{|\mathcal{P}|}\sum_{p}\min_{g}d_{per}(g,p)\Big];
\end{equation}
and the temporal consistency $\mathrm{TC}=\frac1T\sum_j\mathrm{IoU}(\hat M(t_j),\hat M(t_{j-1}))$ with $\mathrm{TCE}=|\mathrm{TC}_{pred}-\mathrm{TC}_{gt}|$. Latency is the mean time to segment one $128\times128$ frame with batch 1, measured after all trainings with a dedicated GPU; memory is the peak GPU memory. Since frames of a trajectory are correlated, 95\% confidence intervals use a bootstrap over test trajectories \citep{efron1993} (2,000 resamples, stratified by regime), and the effect of a technique is the paired difference $\Delta_i=\mathrm{Dice}_i(\text{technique})-\mathrm{Dice}_i(\text{architecture alone})$, bootstrapped in the same way.

%% ===================================================================
\section{Materials and methods}
\label{sec:methods}

\subsection{Data and input resolution}
\label{sec:data}

The protocol limits each subset to 25~GB, prioritising diversity of trajectories and frames. The HDF5 files store each field as a contiguous uncompressed array $[n_{traj},1001,128,128]$, so any frame can be read with an HTTP range request. We used all 1,200 trajectories of the official split files and 21 snapshots spread over the whole trajectory ($n=0,50,\dots,1000$): 960/120/120 trajectories and 20,160/2,520/2,520 frames, 3.3~GB, with no trajectory shared across splits. Using every trajectory keeps all available diversity of initial conditions; the stride of 50 snapshots (500 time units) keeps training cost comparable to a denser sampling of fewer trajectories. Trajectory and frame identifiers, the dataset revision and the final size are recorded in a manifest (Table~\ref{tab:data}). The native $128\times128$ grid (aspect ratio 1:1) is already at the protocol cap, so the field was used without resampling, cropping or padding; the protocol's two-resolution pilot does not apply, since a higher level would require upscaling.

\inputtable{tab_data}

\subsection{Ground truth}
\label{sec:gt}

For each trajectory $i$, the Otsu threshold \citep{otsu1979,sezgin2004} of $B_i$ at the reference snapshot $n_{ref}=500$, $\tau_i=\arg\max_\tau\omega_0\omega_1(\mu_0-\mu_1)^2$, defines the mask of every instant, $M_i(x,y,n)=\ind[B_i(x,y,n)>\tau_i]$, so that the masks record how structures nucleate, grow, split, migrate and vanish. If the reference frame has no contrast ($\max B-\min B\le0.05$), the pattern has collapsed and the median valid threshold of the regime in the training split is used; this happened in 70 of 160 training trajectories of \emph{gliders} and 4 of 160 of \emph{spirals}, plus 12 validation and 9 test trajectories. The ground truth was frozen before training. Networks receive only $A$, never $B$; a learning-free reference, per-frame Otsu thresholding of $A$, $\hat M=\ind[A<\tau_A]$, quantifies what they add to the physical anti-correlation.

\subsection{Configurations}

Table~\ref{tab:methods} lists the methods. The protocol requires a baseline with the chosen standard normalisation and no spatial filter; for Gray--Scott, the standard normalisation is the model's own non-dimensionalisation (concentrations relative to the feed concentration of $A$), so the baseline P0 is the raw dimensionless field and the filters P6--P10 are applied to it, which separates the effect of normalisations (P1--P5 vs.\ P0) from that of filters (P6--P10 vs.\ P0). The percentiles of P5 were inspected before use: $p_1=0.252$ and $p_{99}=1.000$ in training, so clipping saturates the low concentrations inside structures and the excess above 1 present only in initial conditions. Post-processing parameters (threshold, structuring-element radius, minimum area and temporal window) were selected on validation by mean Dice and frozen before testing.

\inputtable{tab_methods}

All backbones started from RGB pre-trained weights, with the first convolution adapted to one channel by summing its kernels over colour. Following the protocol, all networks were trained for exactly 20 epochs (one epoch = 20,160 frames, batch 32) without early stopping, keeping the checkpoint with the best validation Dice and recording its epoch. Optimisation used AdamW \citep{loshchilov2019adamw} (weight decay $10^{-4}$; learning rates $3\times10^{-4}$, $2\times10^{-4}$ and $10^{-4}$ for U-Net, SegFormer and Mask2Former), one warm-up epoch and cosine decay \citep{loshchilov2017sgdr}, and bfloat16. U-Net and SegFormer minimise binary cross-entropy plus soft Dice \citep{milletari2016vnet}. Circular shifts, $90^\circ$ rotations and reflections, exact symmetries of the isotropic periodic system, were used as augmentation. Mask2Former was trained with its set-prediction loss on the final decoder layer only, without deep supervision; architecture and inference are unchanged. The reason is cost: the Hungarian matching is solved per sample and per layer on the CPU, and a training iteration took 2.67~s with deep supervision and 0.79~s without it on the GPU used.

\subsection{Conditions, numbering and environment}

For each architecture, the four conditions were selected on validation by mean Dice---architecture alone (P0+Q0), best P1--P10 with Q0, P0 with best Q1--Q5, and best pair---and evaluated once on test. Following the class convention, E001--E150 are P1--P10 $\times$ architecture $\times$ Q1--Q5 and E151--E198 the baseline rows. Training ran on an NVIDIA RTX A5000 (24~GB) in a container with about seven CPU cores, with up to five trainings in parallel; the full run took about 11.5~h. The code uses PyTorch \citep{paszke2019pytorch}, Hugging Face Transformers \citep{wolf2020transformers}, Segmentation Models PyTorch \citep{iakubovskii2019smp} and scikit-image \citep{vanderwalt2014skimage}, seed 42.

%% ===================================================================
\section{Results}
\label{sec:results}

\subsection{Visual tables}

Figs.~\ref{fig:vt_spirals}--\ref{fig:vt_gliders} show, for each architecture, the input, the pre-processed input, the ground truth and the masks of the four conditions at $n=0$, $50$ (first sampled snapshot after the initial condition), $500$ (reference frame) and $1000$; all regimes are in \ref{app:visual}. At $n=0$, configurations without post-processing miss most of the ground truth, whereas configurations with a lowered threshold, and Mask2Former with percentile clipping, segment almost the whole domain. From $n=50$ on, U-Net follows the structures with few errors, whereas SegFormer and Mask2Former concentrate errors at the boundaries.

\begin{figure}[p]
\centering
\fig{visual_table_spirals}
\caption{Visual table, spirals regime: input $A$, pre-processed input, ground truth and masks of the four conditions at $n=0, 50, 500, 1000$. Dark grey: true positive; orange: false positive; blue: false negative.}
\label{fig:vt_spirals}
\end{figure}

\begin{figure}[p]
\centering
\fig{visual_table_gliders}
\caption{Visual table, gliders regime (moving and splitting spots). Same layout as Fig.~\ref{fig:vt_spirals}.}
\label{fig:vt_gliders}
\end{figure}

\subsection{Architectures and conditions (RQ1)}

\inputtable{tab_main}
\inputtable{tab_classes}

U-Net reached a Dice of 0.953 alone and 0.962 with the selected pre- + post-processing (P2+Q1); SegFormer 0.892 alone and 0.892--0.893 with processing; Mask2Former 0.768 alone and 0.784 with P5+Q2 (Table~\ref{tab:main}). The learning-free reference reached 0.888, with a confidence interval overlapping that of SegFormer. Over all snapshots, U-Net exceeds the reference by 6.5--7.4 points, SegFormer is not distinguishable from it, and Mask2Former stays 10.5--12.0 points below it. The reference has recall 0.967 and precision 0.833: the wider depletion halo of $A$ produces false positives at the boundaries. The per-class metrics (Table~\ref{tab:classes}) show that threshold optimisation raised the structure Dice of U-Net but lowered its background IoU from 0.986 to 0.952 and its mIoU from 0.967 to 0.950, because the selected threshold (0.2) enlarges the region classified as structure.

\subsection{Effect of pre-processing (RQ2) and post-processing (RQ3)}

\inputtable{tab_delta}

For U-Net, pre-processing effects ranged from $-0.64$ (median filtering) to $+0.17$ points (min--max); for SegFormer, all were below 0.07 points in magnitude; for Mask2Former, percentile clipping raised Dice by 1.37 points and the others by 0.00--0.55 (Table~\ref{tab:delta}). Threshold optimisation raised the Dice of U-Net by 0.84 points, with no distinguishable effect on the Transformers. Opening lowered it by 1.28 (U-Net) and 0.25 points (SegFormer); closing changed it by $-0.18$ and $+0.10$ points; connected-component filtering had effects below 0.005 points, since the minimum areas selected on validation (2--8 px) removed almost nothing. Temporal consistency filtering lowered Dice by 11.29, 10.58 and 8.63 points; for U-Net the loss was 31.43 points in \emph{gliders} and 28.95 in \emph{spirals}, whose structures move, against 0.51 in \emph{bubbles} and 0.33 in \emph{worms}. Fig.~\ref{fig:heat} shows all combinations.

\begin{figure}[t]
\centering
\fig{heatmap_dice}
\caption{Test Dice of every combination of pre-processing (rows) and post-processing (columns); P0 and Q0 are the baselines.}
\label{fig:heat}
\end{figure}

\subsection{Sensitivity to the initial snapshot}
\label{sec:sensitivity}

\inputtable{tab_sensitivity}

At $n=0$, one of 21 snapshots, the initial condition is random. Architectures without post-processing reached Dice 0.100--0.120 there, U-Net predicting about 11\% of the domain against 25.3\% in the ground truth. With the 0.2 threshold, U-Net segmented 95.2\% of the domain at $n=0$ and the Dice of that frame rose to 0.353, while over $n\geq50$ the P2+Q1 configuration reached 0.992, below the 0.995 of U-Net alone (Table~\ref{tab:sensitivity}). Recomputed without the initial snapshot, the paired effect of threshold optimisation changed sign: $-0.28$ points for U-Net and $-1.21$ for SegFormer, both distinguishable from zero. The effect of percentile clipping on Mask2Former dropped from 1.37 to 0.18 points, since that configuration also predicted 99.2\% of the domain at $n=0$. The effects of opening and temporal filtering kept their sign and magnitude. Without the initial snapshot, SegFormer reached 0.931, above the 0.921 of the reference.

\subsection{Temporal evolution (RQ4), regimes and cost}

\inputtable{tab_temporal}

At $n=50$, U-Net alone already reached 0.994, SegFormer 0.952 and Mask2Former 0.829, and from $n=100$ on the means were 0.995, 0.930 and 0.800 (Fig.~\ref{fig:temporal}); the reference reached 0.869 and 0.923. U-Net alone had the lowest area-variation error (0.76\% of the domain per sampled step) and TCE 0.010 (ground-truth TC 0.689); its final configuration lowered CDE from 2.68 to 1.42 px but raised AVE to 3.58\%, an effect of the area jump between the over-segmented initial frame and the following ones. SegFormer and Mask2Former kept CDE of 7.1--10.0 px.

\begin{figure}[t]
\centering
\fig{temporal_dice_area}
\caption{Left: mean test Dice per snapshot for the final configuration of each architecture (solid), the architecture alone (dotted) and Otsu on $A$ (dashed). Right: mean segmented area.}
\label{fig:temporal}
\end{figure}

\inputtable{tab_regime}
\inputtable{tab_cost}

U-Net alone kept Dice between 0.943 and 0.979 in all regimes; SegFormer had its lowest values in \emph{spots} (0.833), \emph{maze} (0.865) and \emph{worms} (0.868), below the 0.906, 0.938 and 0.896 of the reference, and Mask2Former in \emph{spots} (0.667) (Table~\ref{tab:regime}). U-Net had the lowest latency (5.58~ms) and Mask2Former the highest (38.63~ms) (Table~\ref{tab:cost}). The best epoch of SegFormer was 19 or 20 for every input, i.e.\ its validation Dice was still increasing at the end of the budget, whereas that of Mask2Former was between epochs 1 and 9, its validation Dice staying between 0.771 and 0.776 for the reference input (Fig.~\ref{fig:training}).

\begin{figure}[t]
\centering
\fig[0.55\linewidth]{training_curves}
\caption{Validation Dice over the 20 training epochs (reference input P0).}
\label{fig:training}
\end{figure}

%% ===================================================================
\section{Discussion}
\label{sec:discussion}

The ranking of the architectures was determined by a design property, the output resolution, more than by the techniques around the network. U-Net predicts at full resolution with skip connections from the highest-resolution stage, whereas SegFormer and Mask2Former predict on a $32\times32$ grid for a $128\times128$ input; with structures a few pixels wide, one output cell covers a large part of a structure and the boundary depends on the interpolation of the logits. This is consistent with the drop of SegFormer in the regimes with the shortest wavelength, where it fell below thresholding of $A$, although the design does not isolate output resolution from backbone and pre-training differences.

Pre-processing had small effects because the simulated fields are noise-free: denoising filters have little to remove, smoothing filters remove signal---the Gaussian filter attenuates the harmonics that make boundaries sharp and the median filter rounds corners and erodes one-pixel necks---and affine normalisations are largely absorbed by the normalisation layers of the networks. The losses of U-Net, the only architecture able to exploit fine detail, are consistent with this.

The sensitivity to the initial snapshot is the result with the largest methodological consequence. At $n=0$ the initial condition is random and $A$ does not yet reflect the structure of $B$, so no configuration segmented that frame informatively. Still, one frame out of 21 was enough to decide the selection on validation: a low threshold makes the network segment almost the whole domain at $n=0$, which increases the Dice of that frame without any physical information, and the same mechanism explains most of the gain of percentile clipping on Mask2Former. Over the frames with a pattern, threshold optimisation lowered the Dice of U-Net and SegFormer and lowered the background IoU, showing that the mean structure Dice, used as the sole selection criterion, favoured over-segmenting an uninformative frame. Choosing a configuration to track the evolution of the structures should therefore consider per-class and temporal metrics and exclude or down-weight the initial snapshot when the initial condition is random.

Post-processing effects followed the scale of the operations relative to the structures and to the temporal sampling. Opening with a radius-1 disk removes structures less than three pixels wide, such as nascent spots and the thin walls of \emph{bubbles}. Connected-component filtering did not change the results because the networks, applied to smooth fields, rarely produce isolated specks, so the hypothesis of point noise that justifies the method did not hold in this dataset. Temporal filtering combines snapshots 50 steps (500 time units) apart; where the pattern moves, the structures change position within that interval and the median produces lagged masks that erase real physical changes---the risk the protocol itself warns about.

Mask2Former should be read with caution: its validation Dice plateaued within the first epochs while the training loss kept decreasing, which is consistent with two hypotheses the design does not separate---the low output resolution and the training without deep supervision, part of the original recipe \citep{cheng2022mask2former}. Its mask-classification formulation was also designed for objects and stuff regions of natural scenes, whereas here one class is spread over tens of disconnected components. The SegFormer best epoch at the end of the budget suggests that the 20-epoch limit of the protocol may have constrained it.

\paragraph{Limitations} One seed per training, so confidence intervals capture variability across test trajectories but not across trainings; Mask2Former without deep supervision; architectures that differ simultaneously in output resolution, backbone and pre-training; selection by mean Dice sensitive to the initial snapshot; training times affected by GPU sharing (latencies were measured with a dedicated GPU); a threshold-based ground truth that measures the recovery of a physically defined segmentation, not the discovery of unknown structures; a stride of 50 snapshots that limits the temporal resolution of the evolution metrics; and one physical system, one GPU and frame-wise networks.

%% ===================================================================
\section{Conclusions}
\label{sec:conclusions}

Under the class protocol, with all 1,200 trajectories of the Gray--Scott system of The Well, U-Net reached a Dice of 0.953 without pre- or post-processing, above the learning-free Otsu reference on $A$ (0.888), whereas SegFormer was not distinguishable from it (0.892) and Mask2Former stayed below it (0.768) (RQ1). Pre-processing changed Dice by at most 1.37 points, the largest effect being driven by the initial snapshot (RQ2). Threshold optimisation raised the mean Dice of U-Net by 0.84 points but lowered it once the initial snapshot was excluded, and temporal consistency filtering lowered the Dice of all architectures by 8.6--11.3 points, with losses of about 30 points in the regimes whose structures move (RQ3). U-Net followed the evolution of the structures from the first sampled snapshot, with a mean Dice of 0.995 from $n=100$ on, but no configuration segmented the random initial snapshot (RQ4). In this dataset, the output resolution of the architecture weighed more than the techniques applied before and after it, and selecting configurations by mean Dice was sensitive to a single frame without physical information. Further work should repeat the trainings with more seeds, train Mask2Former with deep supervision, vary the output resolution of the Transformer architectures, define the selection criterion with the initial snapshot excluded or down-weighted, and merge these results with those of the other physical systems of the class.

\section*{Data and code availability}
The data are publicly available from The Well (\url{https://polymathic-ai.org/the_well/}). The code that downloads the subset, builds the ground truth, trains the models and generates every table and figure is available from the author.

\section*{Declaration of competing interest}
The author declares no known competing financial interests or personal relationships that could have appeared to influence the work reported in this paper.

\bibliographystyle{elsarticle-harv}
\bibliography{references}

\clearpage
\appendix
\section{Method codes}
\label{app:methods}
See Table~\ref{tab:methods}.

\section{Sequential table of all combinations}
\label{app:sequential}
\inputtable{tab_sequential}

\section{Visual tables of the remaining regimes}
\label{app:visual}
\begin{figure}[h!]
\centering
\fig{visual_table_bubbles}
\caption{Visual table, bubbles regime. Same layout as Fig.~\ref{fig:vt_spirals}.}
\end{figure}
\begin{figure}[h!]
\centering
\fig{visual_table_maze}
\caption{Visual table, maze regime.}
\end{figure}
\begin{figure}[h!]
\centering
\fig{visual_table_worms}
\caption{Visual table, worms regime.}
\end{figure}
\begin{figure}[h!]
\centering
\fig{visual_table_spots}
\caption{Visual table, spots regime.}
\end{figure}
\begin{figure}[h!]
\centering
\fig{fluid_by_time}
\caption{Concentration of $B$ for one test trajectory per regime, $n=0,\dots,1000$.}
\end{figure}

\end{document}
